% !TEX root = chapter-standalone.tex

\devel{
\section{A survey of examples}\label{sec:survey-examples}

\todotext{In the section we survey a broad range of examples of (mathematically formulated) structures that compose serially, in order to illustrate the general pattern and motivate the definition of semicategories and categories as a unifying concept.}

In the following we aim to illustrate an overall pattern; we do so by listing a number of examples of series composition for various different mathematical structures encountered in this book thus far. And before listing the examples, let us briefly and informally describe the pattern. In the next section we will then encode this pattern in the formal mathematical definition of a category. 

Loosely speaking, categories describe ``worlds'' which are build of the following constituents: there are things (generically called ``objects''), there are relationships between the things (generically called ``morphisms''), and there is a notion of series composition for those relationships (generically called ``composition''). Furthermore, for each thing (object) there is a special self-relationship called an ``identity morphism''. 

These four constituents are subject to two conditions: that composition must be associative, and that identity morphisms must be neutral for composition. The selected examples below give a first sense of what this may look like. 


}